% ============================================================
%   CHAPITRE : Loi binomiale
%   Niveau   : Terminale S
%   Mazava Maths
% ============================================================
\chapter{Loi binomiale}

\begin{objectifs}
	\begin{itemize}[label=\textbullet]
		\item Calculer un coefficient binomial et utiliser le triangle de Pascal.
		\item Reconnaître une situation relevant de la loi binomiale.
		\item Calculer une probabilité, l'espérance et la variance d'une loi binomiale.
	\end{itemize}
\end{objectifs}


% ============================================================
\section{Coefficients binomiaux}
% ============================================================

\begin{definition}[Coefficient binomial]
	Dans un schéma de Bernoulli à $n$ épreuves, le \textbf{coefficient binomial} $\dbinom{n}{k}$ (noté aussi $C_n^k$) est le nombre de chemins de l'arbre pondéré conduisant à exactement $k$ succès. On a :
	\[
	\binom{n}{k} = \frac{n!}{k!\,(n-k)!}, \qquad 0\leqslant k\leqslant n.
	\]
\end{definition}

\begin{propriete}{Propriétés des coefficients binomiaux}{}
	\begin{itemize}[label=\textbullet]
		\item $\dbinom{n}{0} = \dbinom{n}{n} = 1$.
		\item \textbf{Symétrie} : $\dbinom{n}{k} = \dbinom{n}{n-k}$.
		\item \textbf{Relation de Pascal} : $\dbinom{n-1}{k-1} + \dbinom{n-1}{k} = \dbinom{n}{k}$.
	\end{itemize}
\end{propriete}

La relation de Pascal permet de construire, de proche en proche, le \textbf{triangle de Pascal} : chaque coefficient s'obtient en ajoutant les deux coefficients situés juste au-dessus de lui.

\begin{center}
	\begin{tikzpicture}[font=\small]
		\node at (-3.4,0)    {$n=0$};
		\node at (-3.4,-0.9) {$n=1$};
		\node at (-3.4,-1.8) {$n=2$};
		\node at (-3.4,-2.7) {$n=3$};
		\node at (-3.4,-3.6) {$n=4$};
		\node at (-3.4,-4.5) {$n=5$};

		\node at (0,0) {$1$};

		\node at (-0.5,-0.9) {$1$}; \node at (0.5,-0.9) {$1$};

		\node at (-1,-1.8) {$1$}; \node at (0,-1.8) {$2$}; \node at (1,-1.8) {$1$};

		\node at (-1.5,-2.7) {$1$}; \node at (-0.5,-2.7) {$3$}; \node at (0.5,-2.7) {$3$}; \node at (1.5,-2.7) {$1$};

		\node at (-2,-3.6) {$1$}; \node at (-1,-3.6) {$4$}; \node at (0,-3.6) {$6$}; \node at (1,-3.6) {$4$}; \node at (2,-3.6) {$1$};

		\node at (-2.5,-4.5) {$1$}; \node at (-1.5,-4.5) {$5$}; \node at (-0.5,-4.5) {$10$}; \node at (0.5,-4.5) {$10$}; \node at (1.5,-4.5) {$5$}; \node at (2.5,-4.5) {$1$};
	\end{tikzpicture}
\end{center}

\begin{astucebac}
	Pour lire $\dbinom{n}{k}$ dans le triangle : $n$ désigne le numéro de la ligne (en partant de $0$) et $k$ le numéro de la colonne (en partant de $0$, depuis la gauche).
\end{astucebac}


% ============================================================
\section{Formule du binôme de Newton}
% ============================================================

\begin{theoreme}[Formule du binôme de Newton]
	Pour tous nombres réels $a$, $b$ et tout entier naturel $n$ :
	\[
	(a+b)^n = \sum_{k=0}^{n} \binom{n}{k} a^k b^{n-k}.
	\]
\end{theoreme}

\begin{propriete}{Somme des coefficients binomiaux}{}
	Pour tout entier naturel $n$ :
	\[
	\sum_{k=0}^{n} \binom{n}{k} = 2^n.
	\]
\end{propriete}

\begin{methode}
	\textbf{Démonstration :} il suffit d'appliquer la formule du binôme avec $a=b=1$ :
	\[
	2^n = (1+1)^n = \sum_{k=0}^{n} \binom{n}{k}\,1^k\,1^{n-k} = \sum_{k=0}^{n} \binom{n}{k}.
	\]
\end{methode}

\begin{exemple}
	Pour $n=4$, la ligne du triangle de Pascal donne $1+4+6+4+1 = 16 = 2^4$.
\end{exemple}


% ============================================================
\section{Loi binomiale}
% ============================================================

\begin{definition}[Loi binomiale]
	On considère un schéma de Bernoulli de paramètres $n$ et $p$ (répétition de $n$ épreuves de Bernoulli identiques et indépendantes, de probabilité de succès $p$). Soit $X$ le nombre de succès obtenus.

	On dit que $X$ suit la \textbf{loi binomiale} de paramètres $n$ et $p$, notée $\mathscr{B}(n,p)$, et :
	\[
	\Prob(X=k) = \binom{n}{k}\,p^k\,(1-p)^{n-k}, \qquad k\in\{0,1,\ldots,n\}.
	\]
\end{definition}

\begin{methode}
	\textbf{Reconnaître une loi binomiale :} vérifier que la situation consiste à répéter $n$ fois, de façon \textbf{identique et indépendante}, une épreuve à \textbf{deux issues} (succès/échec) de probabilité de succès $p$ constante, et que $X$ compte le \textbf{nombre de succès}.
\end{methode}

\begin{exemple}
	Un QCM comporte $8$ questions indépendantes, chacune avec $4$ réponses possibles dont une seule est correcte. Un candidat répond au hasard à chaque question. Soit $X$ le nombre de bonnes réponses.

	Chaque question est une épreuve de Bernoulli de paramètre $p=\dfrac14$ (succès = bonne réponse), répétée $8$ fois de façon indépendante : $X$ suit la loi $\mathscr{B}\!\left(8\,;\dfrac14\right)$.

	Par exemple : $\Prob(X=2) = \dbinom{8}{2}\left(\dfrac14\right)^2\left(\dfrac34\right)^6 \approx 0,31$.
\end{exemple}

\begin{center}
	\begin{tikzpicture}[>=stealth, font=\small]
		\draw[->] (-0.6,0) -- (9.5,0) node[right] {$k$};
		\draw[->] (0,-0.3) -- (0,4) node[above] {$\Prob(X=k)$};
		\foreach \x/\h in {0/1.1, 1/2.9, 2/3.4, 3/2.3, 4/1.0, 5/0.25, 6/0.05, 7/0.01, 8/0.002}{
			\draw[very thick, black!70] (\x,0) -- (\x,\h);
			\node[below] at (\x,0) {\scriptsize \x};
		}
	\end{tikzpicture}
\end{center}

\begin{propriete}{Espérance, variance et écart-type de la loi binomiale}{}
	Si $X$ suit la loi $\mathscr{B}(n,p)$, alors :
	\[
	\Esp(X) = np, \qquad \Var(X) = np(1-p), \qquad \sigma(X) = \sqrt{np(1-p)}.
	\]
\end{propriete}

\begin{exemple}
	Pour le QCM précédent : $\Esp(X) = 8\times\dfrac14 = 2$ (en moyenne, un candidat répondant au hasard obtient $2$ bonnes réponses) et $\Var(X) = 8\times\dfrac14\times\dfrac34 = \dfrac32$.
\end{exemple}


% ============================================================
\section{Exercices résolus}
% ============================================================

\exerciceresolu
\textbf{Calculer une probabilité binomiale}

On lance $5$ fois de suite un dé équilibré. On considère comme succès l'obtention d'un $6$. Soit $X$ le nombre de fois où l'on obtient un $6$.
\begin{enumerate}
	\item Justifier que $X$ suit une loi binomiale et préciser ses paramètres.
	\item Calculer $\Prob(X=2)$, arrondie à $10^{-3}$.
	\item Calculer $\Prob(X\geqslant1)$.
	\item Calculer $\Esp(X)$.
\end{enumerate}

\solution

\textbf{1.} Les $5$ lancers sont identiques et indépendants, chacun avec deux issues (« obtenir un $6$ » ou non) de probabilité $p=\dfrac16$ constante : $X$ suit la loi $\mathscr{B}\!\left(5\,;\dfrac16\right)$.

\textbf{2.} $\Prob(X=2) = \dbinom{5}{2}\left(\dfrac16\right)^2\left(\dfrac56\right)^3 = 10\times\dfrac{1}{36}\times\dfrac{125}{216} \approx 0,161$.

\textbf{3.} $\Prob(X\geqslant1) = 1-\Prob(X=0) = 1-\left(\dfrac56\right)^5 \approx 1-0,402 = 0,598$.

\textbf{4.} $\Esp(X) = 5\times\dfrac16 = \dfrac56$.


% ============================================================
\section{Exercices}
% ============================================================

\exercice{Triangle de Pascal \hfill \facile}

\begin{enumerate}
	\item Compléter les deux lignes suivantes du triangle de Pascal après celle des $1\ 5\ 10\ 10\ 5\ 1$.
	\item Vérifier que la somme des coefficients de la ligne $n=6$ vaut $2^6$.
\end{enumerate}

\exercice{Calculs directs \hfill \facile}

Une variable aléatoire $X$ suit la loi binomiale $\mathscr{B}(6\,;0,3)$.
\begin{enumerate}
	\item Calculer $\Prob(X=0)$, $\Prob(X=3)$ et $\Prob(X=6)$.
	\item Calculer $\Esp(X)$ et $\Var(X)$.
\end{enumerate}

\exercice{Tir à l'arc \hfill \moyen}

Un archer a une probabilité de $0,8$ d'atteindre sa cible à chaque tir, indépendamment des tirs précédents. Il tire $6$ flèches. Soit $X$ le nombre de flèches qui atteignent la cible.
\begin{enumerate}
	\item Justifier que $X$ suit une loi binomiale dont on précisera les paramètres.
	\item Calculer la probabilité qu'il atteigne exactement $5$ fois la cible.
	\item Calculer la probabilité qu'il atteigne la cible au moins une fois.
	\item Calculer $\Esp(X)$ et interpréter le résultat.
\end{enumerate}

\exercice{QCM et loi binomiale \hfill \difficile}

Un QCM comporte $10$ questions indépendantes à $5$ réponses possibles chacune, une seule étant correcte. Un candidat répond entièrement au hasard. Soit $X$ le nombre de bonnes réponses.
\begin{enumerate}
	\item Donner la loi de $X$.
	\item Calculer la probabilité d'obtenir au moins $2$ bonnes réponses.
	\item Un candidat est reçu s'il obtient au moins $6$ bonnes réponses. Calculer la probabilité que ce candidat, répondant au hasard, soit reçu.
	\item Calculer $\Esp(X)$ et $\sigma(X)$.
\end{enumerate}
